\begin{center}
    {\large\bfseries 1. Следует или не следует?}
\end{center}

\footnotesize

\begin{center}
\fbox{%
\begin{minipage}{0.91\linewidth}
\textbf{Коротко о главном.} Если из условия $A$ обязательно получается вывод $B$,
то говорят: «из $A$ следует $B$» и пишут $A\Rightarrow B$.
Если можно найти хотя бы один контрпример, то следствия нет.
\end{minipage}}
\end{center}

\begin{enumerate}[itemsep=0.55em]

\item \textbf{Верно или неверно:}

\begin{enumerate}[label={},leftmargin=1.7em,itemsep=0.65em]
\item[а)] Если число делится на $18$, то оно делится на $6$;
\item[б)] если квадрат числа делится на $9$, то само число делится на $9$;
\item[в)] если диагонали прямоугольника равны, то это квадрат;
\item[г)] если число больше $13$, то оно не меньше $12$;
\item[д)] если фигуру одним разрезом разделили на два треугольника, то она была четырёхугольником?
\end{enumerate}

\item \textbf{Что из чего следует?}

В каждой строке соедините утверждения стрелкой в верном направлении.
Если верны оба направления, нарисуйте две стрелки. Если ни одного — поставьте знак~$\times$.

\begin{tabularx}{\linewidth}{@{}>{\raggedright\arraybackslash}X c >{\raggedright\arraybackslash}X@{}}
а) оканчивается цифрой $0$ & \rule{1cm}{0.4pt} & делится на $10$; \\[0.4em]
б) делится на $15$ & \rule{1cm}{0.4pt} & делится на $5$; \\[0.4em]
в) квадрат & \rule{1cm}{0.4pt} & ромб; \\[0.4em]
г) $ab$ делится на $6$ & \rule{1cm}{0.4pt} & $a$ или $b$ делится на $6$; \\[0.4em]
д) $n^2$ делится на $12$ & \rule{1cm}{0.4pt} & $n$ делится на $6$.
\end{tabularx}

Здесь $a$, $b$ и $n$ — натуральные числа.

\newpage

\item \textbf{Карта следствий.}

Даны шесть свойств натурального числа. Проведите стрелки так, чтобы из схемы
можно было увидеть все верные следствия. Стрелка идёт от условия к его следствию;
лишних стрелок не проводите.

\begin{center}
\begin{tikzpicture}[
    box/.style={draw,rounded corners=2pt,align=center,minimum width=3.25cm,minimum height=0.9cm,font=\small},
    >=stealth
]
    \node[box] (a) at (0,2.2) {делится на $36$};
    \node[box] (b) at (-2.35,0.7) {делится на $18$};
    \node[box] (c) at (2.35,0.7) {делится на $12$};
    \node[box] (d) at (-2.35,-0.8) {делится на $9$};
    \node[box] (e) at (2.35,-0.8) {делится на $6$};
    \node[box] (f) at (0,-2.3) {делится на $3$};
\end{tikzpicture}
\end{center}

\item \textbf{Выберите все следствия.}

Известно, что натуральное число делится на $12$. Какие утверждения обязательно верны?
Обведите их номера.

\begin{multicols}{2}
\begin{enumerate}[label=\arabic*),itemsep=0.35em]
\item число делится на $2$;
\item число делится на $3$;
\item число делится на $4$;
\item число делится на $5$;
\item число делится на $6$;
\item число делится на $8$;
\item число чётное;
\item сумма его цифр делится на $3$.
\end{enumerate}
\end{multicols}

\item \textbf{Цепочки следствий.}

Продолжите каждую цепочку и запишите общий вывод.

\begin{enumerate}[label={},leftmargin=1.7em,itemsep=0.75em]
\item[а)] Если число делится на $24$, то оно делится на $8$.
Если число делится на $8$, то оно чётное.
Следовательно, если число делится на $24$, то \rule{4.4cm}{0.4pt}.

\item[б)] Если четырёхугольник — квадрат, то он является прямоугольником.
У всякого прямоугольника диагонали равны.
Следовательно, у квадрата \rule{4.5cm}{0.4pt}.

\item[в)] Если сумма цифр числа делится на $9$, то само число делится на $9$.
Если число делится на $9$, то оно делится на $3$.
Следовательно, \rule{7.6cm}{0.4pt}.
\end{enumerate}

\item \textbf{Охота за контрпримером.}\par\vspace{0.15em}

Каждое следствие ниже неверно. Подберите контрпример.

\begin{enumerate}[label={},leftmargin=1.7em,itemsep=0.75em]
\item[а)] Если число делится на $4$, то оно делится на $8$.

Контрпример: \rule{6.8cm}{0.4pt}

\item[б)] Если сумма двух натуральных чисел чётна, то оба слагаемых чётны.

Контрпример: \rule{6.8cm}{0.4pt}

\item[в)] Если у прямоугольника равны диагонали, то он является квадратом.

Контрпример: \rule{6.8cm}{0.4pt}

\item[г)] Если натуральное число больше $10$, то оно делится на $2$ или на $3$.

Контрпример: \rule{6.8cm}{0.4pt}
\end{enumerate}

\item \textbf{Составьте сами.}

Используя одно из условий ниже, составьте одно верное и одно неверное следствие.
Для неверного следствия укажите контрпример.

\begin{center}
\fbox{число делится на $20$}
\qquad
\fbox{треугольник равносторонний}
\qquad
\fbox{$x>100$}
\end{center}

Верное: \rule{8.7cm}{0.4pt}

Неверное: \rule{8.5cm}{0.4pt}

Контрпример: \rule{7.8cm}{0.4pt}

\end{enumerate}
